\section{Split the source into files}
\label{sec:intro}

This section lives in \texttt{sections/introduction.tex}.
\texttt{main.tex} pulls it in with:

\begin{verbatim}
\section{Split the source into files}
\label{sec:intro}

This section lives in \texttt{sections/introduction.tex}.
\texttt{main.tex} pulls it in with:

\begin{verbatim}
\section{Split the source into files}
\label{sec:intro}

This section lives in \texttt{sections/introduction.tex}.
\texttt{main.tex} pulls it in with:

\begin{verbatim}
\section{Split the source into files}
\label{sec:intro}

This section lives in \texttt{sections/introduction.tex}.
\texttt{main.tex} pulls it in with:

\begin{verbatim}
\input{sections/introduction}
\end{verbatim}

That is the usual way to keep a growing paper readable: one file per section,
plus a short root file that only sets up packages and the document skeleton.

\subsection{Shared macros}
Commands such as \texttt{\textbackslash R} and \texttt{\textbackslash norm}
are defined once in \texttt{macros.tex} (loaded near the top of
\texttt{main.tex}). After that, you can write $\R$ or $\norm{x}$ anywhere
without repeating the definition.

\subsection{A labeled equation}
For $x \in \R$, a classic identity is
\begin{equation}
\label{eq:euler}
  e^{i\pi} + 1 = 0.
\end{equation}
The \texttt{\textbackslash label\{eq:euler\}} makes it possible to point back
here later with \texttt{\textbackslash eqref\{eq:euler\}} or
\texttt{\textbackslash ref\{eq:euler\}}. Cross-references update after you
compile again (a second pass is normal).

\subsection{Citations from a \texttt{.bib} file}
A citation like \cite{lamport94} comes from \texttt{references.bib}.
After the first successful compile, run your bibliography tool (BibTeX), then
compile again so the reference list and in-text citations resolve.

\end{verbatim}

That is the usual way to keep a growing paper readable: one file per section,
plus a short root file that only sets up packages and the document skeleton.

\subsection{Shared macros}
Commands such as \texttt{\textbackslash R} and \texttt{\textbackslash norm}
are defined once in \texttt{macros.tex} (loaded near the top of
\texttt{main.tex}). After that, you can write $\R$ or $\norm{x}$ anywhere
without repeating the definition.

\subsection{A labeled equation}
For $x \in \R$, a classic identity is
\begin{equation}
\label{eq:euler}
  e^{i\pi} + 1 = 0.
\end{equation}
The \texttt{\textbackslash label\{eq:euler\}} makes it possible to point back
here later with \texttt{\textbackslash eqref\{eq:euler\}} or
\texttt{\textbackslash ref\{eq:euler\}}. Cross-references update after you
compile again (a second pass is normal).

\subsection{Citations from a \texttt{.bib} file}
A citation like \cite{lamport94} comes from \texttt{references.bib}.
After the first successful compile, run your bibliography tool (BibTeX), then
compile again so the reference list and in-text citations resolve.

\end{verbatim}

That is the usual way to keep a growing paper readable: one file per section,
plus a short root file that only sets up packages and the document skeleton.

\subsection{Shared macros}
Commands such as \texttt{\textbackslash R} and \texttt{\textbackslash norm}
are defined once in \texttt{macros.tex} (loaded near the top of
\texttt{main.tex}). After that, you can write $\R$ or $\norm{x}$ anywhere
without repeating the definition.

\subsection{A labeled equation}
For $x \in \R$, a classic identity is
\begin{equation}
\label{eq:euler}
  e^{i\pi} + 1 = 0.
\end{equation}
The \texttt{\textbackslash label\{eq:euler\}} makes it possible to point back
here later with \texttt{\textbackslash eqref\{eq:euler\}} or
\texttt{\textbackslash ref\{eq:euler\}}. Cross-references update after you
compile again (a second pass is normal).

\subsection{Citations from a \texttt{.bib} file}
A citation like \cite{lamport94} comes from \texttt{references.bib}.
After the first successful compile, run your bibliography tool (BibTeX), then
compile again so the reference list and in-text citations resolve.

\end{verbatim}

That is the usual way to keep a growing paper readable: one file per section,
plus a short root file that only sets up packages and the document skeleton.

\subsection{Shared macros}
Commands such as \texttt{\textbackslash R} and \texttt{\textbackslash norm}
are defined once in \texttt{macros.tex} (loaded near the top of
\texttt{main.tex}). After that, you can write $\R$ or $\norm{x}$ anywhere
without repeating the definition.

\subsection{A labeled equation}
For $x \in \R$, a classic identity is
\begin{equation}
\label{eq:euler}
  e^{i\pi} + 1 = 0.
\end{equation}
The \texttt{\textbackslash label\{eq:euler\}} makes it possible to point back
here later with \texttt{\textbackslash eqref\{eq:euler\}} or
\texttt{\textbackslash ref\{eq:euler\}}. Cross-references update after you
compile again (a second pass is normal).

\subsection{Citations from a \texttt{.bib} file}
A citation like \cite{lamport94} comes from \texttt{references.bib}.
After the first successful compile, run your bibliography tool (BibTeX), then
compile again so the reference list and in-text citations resolve.
